\section{Prospective features and controlled comparisons}
\label{sec:design}

For readability, G-small, G-match, Atom, Rect, Union and Relative denote the
six transferred families. Regret, Gain, Defer and SPO+ denote the additional
development objectives. Table~\ref{tab:implementation-panel} binds these
display names to their unchanged canonical identities, inputs and outputs.

% Methods source map: M04. EXP-033--038, source authority PROSPECTIVE_V1.
\subsection{A prospective schema, with separate cost accounting}
The historical diagnostic schema, \texttt{RETRO\_V4}, supplied the realized
runtime of the action being predicted in two input coordinates. A zero
decision-cost multiplier did not remove these predictive covariates.
Historical rectangle and union geometry also used the target-validity mask.
The prospective study therefore changes runtime availability, geometry, and
descriptor boundaries together; its comparison with the historical workflow
is not a single-factor causal runtime ablation.

The replacement \texttt{PROSPECTIVE\_V1} schema constructs candidate
descriptors from action plans, removes unexecuted-action runtime inputs, and
keeps reference-derived validity in target artifacts. Already incurred base
cost remains observable. Residual predictors receive no candidate-cost
array. Costs enter Equation~\ref{eq:unified-value} separately: each action
type has the median positive feasible runtime from outer-training groups,
normalized by the lower empirical median pooled over those groups' feasible
repair actions. Neither a held-out state's runtime nor its reference label
enters this profile. Full-fit profiles use the authorized VOC development
groups. These are fixed cost estimates rather than predictions of a test
action's eventual execution time.

% Methods source map: M05. EXP-034--035.
\subsection{Capacity and localization}
The first prospective comparison separates a small global residual model,
\model{R0\_SMALL\_P}, from a larger global control,
\model{R0\_CM\_P}. A target-independent integer width search matches the
latter's parameter count as closely as possible to the local
\model{R1\_P} architecture, with smaller width breaking ties.
\model{RECT\_P} and \model{UNION\_P} retain the shared local architecture
with different full-raster partitions. Rect uses a deterministic guillotine
partition with the same \emph{number} of regions as Atom. The implementation
balances successive rectangle splits; despite the historical ``area-matched''
name, it does not match each atom's area or the full area distribution.
Union uses four-connected components of the candidate union \emph{and its
complement}. Both controls cover every raster pixel exactly once. Their
descriptor dimensions are shared with Atom, but descriptor contents are
adapted to each representation. They test matched local model recipes,
rather than changing only a boundary coordinate.


Figure~\ref{fig:representations}e contrasts the two computation paths.
The local network encodes every region, pools the unweighted mean and maximum
of region embeddings into a state context, and combines that context with
the region embedding and an action embedding. A shared sigmoid head predicts
the regional residual for each action; full-raster area fractions aggregate
these outputs. Thus context pooling and final residual aggregation use
different reductions. Global predictors consume the state descriptors and
one action descriptor without a region encoder. Complete dimensions,
activation functions and feature definitions appear in
Appendix~\ref{sec:implementation-details}.

All first-batch models optimize the same state-level Huber objective and use
the same inner-stopping objective. Local target values are available for the
accounting audit but add no auxiliary supervision to this comparison.
Architectures, folds, seeds, optimization, and stopping are frozen together.
V6 uses one full-partition optimization step per epoch with AdamW, learning
rate $5\times10^{-4}$, weight decay $10^{-4}$, at most 120 epochs, and
patience 15 on the seed-specific inner-stop objective. A decrease must
exceed $10^{-8}$ to replace the retained checkpoint. Feature standardization
is fitted on the outer-training groups, including inner-stop groups but
excluding the outer-held-out fold, after prospective action projection.
The comparison thus addresses capacity and localization within a specified
feature family; parameter matching does not make the architectures identical
or establish a universal causal effect of capacity.

% Methods source map: M06. EXP-036; phase2 semantic addendum dominates old prose.
\subsection{Decision objectives and task adaptations}
\model{R2\_P} adds a soft-policy expected-regret objective to absolute and
atom-local supervision. \model{R3\_P} retains those supervised terms and
adds STOP-relative regression, beneficial-sign supervision, and informative
within-state ordering. Their exact frozen losses appear in
Supplement~\ref{sec:accounting-appendix}. These variants evaluate complete
objective families: because their atom auxiliary terms differ from the
state-only first-batch \model{R1\_P}, they are not a one-term loss ablation.

\model{Q2\_P} directly predicts STOP-relative gain. Its scientific STOP
prediction is fixed to zero before metrics or selection; it has no absolute
residual MAE. \model{L2D\_P} learns masked classification of the feasible
oracle action, and \model{SPO\_PLUS\_P} learns action values with the
finite-action SPO+ surrogate. Both are task adaptations with separate
registered models for each cost multiplier, not claims to reproduce the
source methods' original applications. Their decision outputs have no
residual-error interpretation. \model{LL4TTA\_P} denotes the global-loss
adaptation already instantiated by \model{R0\_SMALL\_P}; it contributes
neither an independent checkpoint family nor another replication.

% Methods source map: M07. EXP-037; prospective phase2 F0/F1 authority.
\subsection{Remaining-action fidelity after a fixed probe}
Because $A3$ is unconditional, the observation study compares
$\mathcal F_0$ with $\mathcal F_1=\mathcal F_0\vee\sigma(A3)$ after that
probe has been executed. The matched \model{F0\_P} and \model{F1\_P}
models rank only $\{\operatorname{STOP},A4,A5,A6\}$ with the same
architecture and state-level objective. Only F1 observes the executed A3
mask, status, and already incurred cost. Both exclude future A4--A6 outcomes.
The reported comparison measures remaining-action MAE and DRRE after the
probe. Practical ties in these metrics establish neither a conditional
utility tie nor absence of information value. The sunk A3 cost is not added
as a new term in this comparison, which therefore does not establish the
end-to-end value of acquiring A3. Other information encodings, probes and
sequential policies remain outside this control.

Continuous predictions are averaged across the registered seeds after fold
stitching and before decisions.
